\documentclass[conference]{IEEEtran}

\usepackage{amsmath}
\usepackage{graphicx}
\usepackage{booktabs}
\usepackage{listings}
\usepackage{xcolor}
\usepackage{hyperref}

\lstset{
  basicstyle=\ttfamily\scriptsize,
  columns=fullflexible,
  breaklines=true,
  backgroundcolor=\color{black!4},
  frame=single,
  language=C++
}

\begin{document}

\title{A Verified CPU/GPU Solver for the 3D Heat Equation}

\author{
  \IEEEauthorblockN{Karl Keshavarzi}
  \IEEEauthorblockA{
    Department of Electrical and Computer Engineering\\
    University of Waterloo\\
    \href{https://github.com/karl-kes/3d-heat-solver}{github.com/karl-kes/3d-heat-solver}
  }
}

\maketitle

\begin{abstract}
We present a C++/CUDA solver for the three-dimensional heat equation using an explicit
(forward Euler) finite-difference scheme, with CPU (OpenMP/SIMD) and CUDA GPU backends
compiled from the same source tree and selected at compile time. We verify both
backends independently against a closed-form analytic solution, a conservation law, an
exact boundary-condition invariant, and a grid-convergence study, rather than treating
performance as a proxy for correctness. The verification caught a genuine defect in the
original GPU boundary kernel and revealed an invalid initial convergence-study design;
the corrected study recovers the expected second-order accuracy. We report benchmark
results from $8^3$ to $512^3$ grid cells on stated hardware, with an analytic
bandwidth argument supporting the memory-bound conclusion in place of an unmeasured
assertion, showing a crossover from CPU-favored to GPU-favored execution between
$16^3$ and $32^3$ cells and a speedup plateau around $65$--$68\times$.
\end{abstract}

\section{Introduction}

The heat equation is a standard model problem for explicit time-stepping
finite-difference methods, and a common first project for learning GPU programming: the
update rule is a simple, highly parallel stencil, but a correct implementation still
requires getting boundary conditions, stability, and memory layout right on both a
CPU and a GPU backend. This paper presents and verifies such an implementation, with
an emphasis on verifying correctness rather than only reporting a speedup number. A
simulation that runs fast but computes the wrong answer is not useful; we treat the
performance results in Section~\ref{sec:perf} as meaningful only because
Section~\ref{sec:verification} independently establishes that both backends solve the
stated equation correctly.

The contributions of this work are:
\begin{itemize}
  \item A dual-backend (OpenMP/SIMD CPU, CUDA GPU) solver for the 3D heat equation
    built from one source tree, with the backend selected at compile time.
  \item A verification methodology combining an analytic-solution check, a
    conservation check, an exact boundary-condition invariant, and a grid-convergence
    study (Section~\ref{sec:convergence}, including a corrected mesh-refinement scheme
    after an initial attempt was found to conflate resolution with problem
    difficulty), applied independently to both backends.
  \item Identification and correction of a genuine GPU boundary-kernel defect that
    this verification methodology caught and a coarser check would have missed
    (Section~\ref{sec:bug}).
  \item A performance characterization from $8^3$ to $512^3$ cells on stated hardware,
    with an analytic argument for the memory-bound conclusion rather than an
    unsupported assertion (Section~\ref{sec:perf}).
\end{itemize}

\subsection{Related work}

GPU finite-difference stencil computation is well studied; Micikevicius
\cite{micikevicius2009} describes shared-memory tiling and multi-GPU strategies for
3D finite-difference stencils on CUDA hardware, in the context of seismic modeling.
That work targeted Tesla/Fermi-era GPUs, which had no general-purpose L1 cache for
global memory, making explicit tiling necessary to capture neighbor reuse at all. The
present work targets a single GPU and does not implement shared-memory tiling
(Section~\ref{sec:limitations}); instead it relies on the much larger L1/L2 cache
hierarchy of modern hardware (Ampere, here) to capture some of that reuse implicitly,
consistent with the cache-reuse factor inferred in Section~\ref{sec:perf}.

\section{Mathematical Formulation}

\subsection{Governing equation}

We solve the unsteady heat equation on a cuboid domain,
\begin{equation}
  \frac{\partial u}{\partial t} = \alpha \nabla^2 u,
\end{equation}
where $u(x,y,z,t)$ is temperature and $\alpha$ is the (constant, isotropic) thermal
diffusivity. The domain is discretized on a structured grid of $n_x \times n_y \times
n_z$ interior cells with uniform spacings $\Delta x, \Delta y, \Delta z$, and the
Laplacian is approximated with the standard second-order central-difference 7-point
stencil,
\begin{equation}
  \nabla^2 u_{i,j,k} \approx
    \frac{u_{i-1,j,k} - 2u_{i,j,k} + u_{i+1,j,k}}{\Delta x^2} + \;(\text{similarly in } y, z).
\end{equation}

\subsection{Time integration and stability}

Time advancement uses explicit (forward) Euler,
\begin{equation}
  u_{i,j,k}^{t+1} = u_{i,j,k}^{t} + \alpha\,\Delta t \, \nabla^2 u_{i,j,k}^{t}.
\end{equation}
This scheme is conditionally stable; von Neumann stability analysis
\cite{obrien1951} of the 3D case gives the classical Courant--Friedrichs--Lewy (CFL)
limit \cite{courant1928}
\begin{equation}
  \Delta t \leq \frac{1}{2\alpha\left(\Delta x^{-2} + \Delta y^{-2} + \Delta z^{-2}\right)}.
  \label{eq:cfl}
\end{equation}
The implementation computes $\Delta t$ from Eq.~\eqref{eq:cfl} with an $0.85$ safety
factor, so that changing the grid spacing or diffusivity from the command line cannot
silently push the simulation into an unstable regime.

\subsection{Boundary conditions}

The domain boundary is insulated: zero heat flux normal to each face,
$\partial u/\partial n = 0$. This is implemented with the standard ghost-cell mirroring
technique, $u_{\text{ghost}} = u_{\text{adjacent interior}}$, applied independently on
each of the six faces after every integration step. Because the central-difference
stencil's flux term at the boundary depends only on the difference between the ghost
and adjacent-interior values, this mirroring makes that difference exactly zero, not
merely approximately so.

\section{Implementation}

\subsection{Dual-backend architecture}

Both the CPU and GPU code paths live in the same \texttt{.cu}/\texttt{.cuh} source
files, selected at compile time via \texttt{\#if defined(\_\_CUDACC\_\_)}. When no CUDA
toolchain is available, CMake compiles the same sources as ordinary C++, and the
project builds and runs with a pure CPU backend. This keeps the two implementations of
each kernel textually adjacent, which made it straightforward to compare them
directly when investigating the boundary-condition defect described in
Section~\ref{sec:bug}.

\subsection{Memory layout}

The field is stored in an aligned, struct-of-arrays container (\texttt{AlignedSoA}).
Each dimension is padded up to a multiple of the SIMD width (16, 32, or 64 bytes,
detected from \texttt{\_\_AVX2\_\_}/\texttt{\_\_AVX512F\_\_}) so that CPU stencil loops
can be vectorized without unaligned loads, and so that the same padded layout is used
unmodified on the GPU. Allocation is dispatched per backend: \texttt{cudaMalloc} under
CUDA, \texttt{\_aligned\_malloc} on MSVC, and \texttt{std::aligned\_alloc} elsewhere.

\subsection{CPU parallelization}

The CPU integration kernel parallelizes the two outer loop dimensions with
\texttt{\#pragma omp parallel for collapse(2)} and vectorizes the innermost loop with
\texttt{\#pragma omp simd} \cite{dagum1998}, relying on \texttt{\_\_restrict\_\_}-qualified
pointers and an alignment hint (\texttt{\_\_builtin\_assume\_aligned} on GCC/Clang,
\texttt{\_\_assume} on MSVC) so the compiler can generate aligned vector loads.

\subsection{GPU kernel design}

The integration kernel launches one thread per interior cell with a $8\times8\times8$
thread block \cite{nvidia-cuda}, each thread computing one stencil update directly from global memory.
The boundary condition is applied with three separate kernels, one per pair of
opposing faces, each launched with a 2D thread grid sized to that face rather than to
the full interior volume.

\subsubsection{A correctness defect found by testing}
\label{sec:bug}

The original boundary-condition kernel used a single 3D-launched kernel that derived
all three face mirrors from one shared, interior-only-ranged thread index. Because
that index never took the value $0$ or $n-1$ in any dimension, edge cells where two
boundary faces meet (e.g. $i=0, j=0$, any interior $k$) never received their second
face's mirror update on the GPU, while the CPU implementation, which used a full range
on the cross-axis in each face's loop, did not have this defect. The discrepancy is
invisible to a coarse correctness check: a center-value comparison never samples an
edge, and a total-energy check is not sensitive to a thin line of stale cells. It was
caught by an explicit per-face boundary test (Section~\ref{sec:verification}) and fixed
by splitting the single kernel into three per-face kernels with correctly-sized 2D
launches, mirroring the CPU code's loop ranges exactly. This also removed an efficiency
problem in the original kernel, which launched $O(n^3)$ threads to perform $O(n^2)$
worth of boundary work.

\section{Verification}
\label{sec:verification}

Following standard code-verification practice \cite{oberkampf2002}, we check the
implementation against problems with a known closed-form solution rather than relying
on benchmark performance as a proxy for correctness. Because the CPU and GPU backends
are selected at compile time, a single binary cannot
run both, so direct cross-binary comparison is not a natural fit. Instead, the same
test suite runs independently against the analytic solution in both build
configurations; if both backends agree with the known-correct answer within a
tolerance $\epsilon$, the triangle inequality bounds their mutual disagreement by at
most $2\epsilon$, without needing a dump-and-diff mechanism. Four checks are used:
three on a small ($32^3$) grid with an isotropic
Gaussian initial condition, and a fourth (Section~\ref{sec:convergence}) sweeping
multiple grid sizes.

\subsection{Analytic solution check}

The initial condition is an isotropic Gaussian, $u(\mathbf{r}, 0) = A_0 \exp(-r^2 /
2\sigma_0^2)$. For the heat equation on an effectively unbounded domain, this remains
Gaussian for all $t$, with
\begin{equation}
  \sigma^2(t) = \sigma_0^2 + 2\alpha t, \qquad
  u(\mathbf{r}, t) = A_0 \left(\frac{\sigma_0^2}{\sigma^2(t)}\right)^{3/2}
    e^{-r^2 / 2\sigma^2(t)}.
\end{equation}
The test runs the simulation for a fixed number of steps and compares the value at the
grid point nearest the center against this closed form, within a $5\%$ relative
tolerance accounting for discretization error and the small residual influence of the
(finite, insulated) domain boundary.

\subsection{Conservation of energy}

With zero flux at every boundary and no source term, $\int u \, dV$ is exactly
conserved: $\frac{d}{dt}\int u\,dV = \alpha \oint \nabla u \cdot \mathbf{n} \, dA = 0$.
The test compares the discrete sum $\sum_{i,j,k} u_{i,j,k}$ at $t=0$ against the same
sum after the run, within a $2\%$ tolerance.

\subsection{Exact boundary invariant}

Because the boundary kernel implements $u_{\text{ghost}} = u_{\text{adjacent interior}}$
as a literal copy with no arithmetic, this is checked for \emph{exact} equality
(no tolerance) at every one of the six faces after the run. This is the test that
caught the defect in Section~\ref{sec:bug}.

\subsection{Grid convergence}
\label{sec:convergence}

The checks above confirm the simulation is close to the analytic solution at one
resolution; they do not confirm the scheme achieves its \emph{designed} second-order
spatial accuracy as the mesh is refined. We measured the global $L_2$ relative error
against the analytic solution over $n \in \{16, 32, 64, 128, 256\}$, evaluated at
every grid point rather than only the center (a single point can show artificially
elevated order from symmetry-driven error cancellation).

A first attempt held $\Delta x = 1$ fixed and varied only $n$; this looked like mesh
refinement but wasn't: the initial condition's width $\sigma_0 = 0.1n$ grew with $n$
while the simulated time stayed fixed, so the Fourier number
$\mathrm{Fo} = \alpha t/\sigma_0^2$ shrank as $1/n^2$ -- the blob diffused less at
each "finer" resolution, a progressively easier target rather than a real accuracy
test, producing a false order of $3.5$--$5$. Fixing the physical domain and
$\sigma_0$ instead, and shrinking $\Delta x = L/n$, holds $\mathrm{Fo}$ approximately
constant ($0.14$--$0.17$) across the sweep and gives the textbook order
$\approx 2.04$ over $n=16$--$128$ (Table~\ref{tab:convergence}, Fig.~\ref{fig:convergence}).
We report the fitted order over this range, where the asymptotic second-order regime
is observed; the order degrades to $1.20$ between $n=128$ and $256$. The CFL limit
ties $\Delta t \propto \Delta x^2$, so this refinement quadruples the step count at
each level ($161$ steps at $n=128$, $641$ at $n=256$); accumulating that many
single-precision increments onto an $O(1)$ field is a plausible source of the
degradation, though we have not isolated it from the first-order-in-time forward
Euler term with a controlled (e.g. double-precision) comparison.

Across all refinement levels, CPU and GPU solutions agreed to displayed precision at
every grid point.

\begin{table}[t]
\centering
\caption{Global $L_2$ relative error, fixed-domain refinement.}
\label{tab:convergence}
\begin{tabular}{@{}cc@{}}
\toprule
$n$ & $L_2$ relative error \\
\midrule
16  & $8.30 \times 10^{-3}$ \\
32  & $1.79 \times 10^{-3}$ \\
64  & $4.29 \times 10^{-4}$ \\
128 & $1.19 \times 10^{-4}$ \\
256 & $5.20 \times 10^{-5}$ \\
\bottomrule
\end{tabular}
\end{table}

\begin{figure}[t]
\centering
\includegraphics[width=\linewidth]{figures/convergence.pdf}
\caption{Global $L_2$ relative error vs. grid size for the invalid self-similar
sweep and the corrected fixed-domain sweep, with the fitted order for the latter.}
\label{fig:convergence}
\end{figure}

\subsection{Continuous integration}

A GitHub Actions pipeline runs the full test suite against the CPU backend on every
push, since hosted runners have no GPU available to execute CUDA kernels. The GPU
backend is compiled in CI to catch build breakage (\texttt{nvcc} does not require a
physical device to compile, only to run), but its tests are run manually against real
hardware.

\section{Performance Evaluation}
\label{sec:perf}

\subsection{Experimental setup}

All results were measured on: CPU, AMD Ryzen 5 8500G (6 cores / 12 threads); GPU,
NVIDIA GeForce RTX 3070, 8\,GB GDDR6 (256-bit bus, spec-sheet peak memory bandwidth
448\,GB/s, peak FP32 throughput approximately 20.3\,TFLOPS); 16\,GB system RAM.

\begin{table}[t]
\centering
\caption{Runtime for 1000 forward-Euler steps, averaged over 5 of 6 runs (first
discarded as warm-up).}
\label{tab:perf}
\begin{tabular}{@{}lrrr@{}}
\toprule
Grid size & CPU (ms) & GPU (ms) & Speedup \\
\midrule
$8^3$   & 4.14    & 32.2 & 0.13$\times$ \\
$16^3$  & 13.8    & 22.7 & 0.61$\times$ \\
$32^3$  & 83.7    & 26.7 & 3.13$\times$ \\
$64^3$  & 535     & 31.8 & 16.8$\times$ \\
$128^3$ & 3{,}805 & 62.0 & 61.4$\times$ \\
$256^3$ & 29{,}856 & 446 & 67.0$\times$ \\
$512^3$ & 225{,}321 & 3{,}336 & 67.6$\times$ \\
\bottomrule
\end{tabular}
\end{table}

\begin{figure}[t]
\centering
\includegraphics[width=\linewidth]{figures/runtime.pdf}
\caption{Runtime vs. grid size for both backends (log-log).}
\label{fig:runtime}
\end{figure}

\begin{figure}[t]
\centering
\includegraphics[width=\linewidth]{figures/speedup.pdf}
\caption{Speedup (CPU time / GPU time) vs. grid size. The dashed line marks parity.}
\label{fig:speedup}
\end{figure}

Table~\ref{tab:perf} and Figs.~\ref{fig:runtime}--\ref{fig:speedup} summarize benchmark
results measured on the hardware described above. At $8^3$ and $16^3$
cells, fixed CUDA kernel-launch overhead (four launches per step: one integration
kernel and three boundary kernels) exceeds the actual compute, and the CPU backend is
faster. The crossover occurs between $16^3$ and $32^3$ cells. Past the crossover,
speedup grows quickly and then plateaus around $65$--$68\times$ from $128^3$ cells
onward.

\subsection{Bandwidth argument}

We support the memory-bandwidth-bound hypothesis analytically rather than with an
unmeasured assertion. The integration kernel's arithmetic intensity is fixed by the
stencil: $16$ floating-point operations and up to $8$ float reads/writes ($32$ bytes)
per cell update in the naive (uncached) case, giving $\mathrm{AI} \approx 0.5$
FLOP/byte. The RTX 3070's roofline ridge point is at
$20.3\,\text{TFLOPS} / 448\,\text{GB/s} \approx 45$ FLOP/byte, two orders of magnitude
above our $\mathrm{AI}$, which places this kernel firmly in the bandwidth-bound region
regardless of the achieved fraction of peak bandwidth \cite{williams2009}.

We can also bound the cache-reuse factor without a profiler. At $512^3$ cells, the
naive $32$ bytes/cell figure implies $4.295 \times 10^9$ bytes moved per integration
step; dividing by this kernel's share of the measured per-step time
($3.336\,\text{ms}$, from $3{,}336\,\text{ms}/1000$ steps in Table~\ref{tab:perf})
gives an implied bandwidth of approximately $1{,}287\,\text{GB/s}$ -- $2.9\times$ the
RTX 3070's $448\,\text{GB/s}$ spec-sheet peak. Since achieved bandwidth cannot exceed
peak, this is direct evidence (not an assumption) that L1/L2 cache reuse across the
7-point stencil's overlapping neighbor reads reduces actual DRAM traffic to well under
a third of the naive estimate. Direct profiler measurement of the achieved-bandwidth
percentage is left as future work.

\section{Discussion and Limitations}
\label{sec:limitations}

The solver targets a single GPU and a single explicit time-stepping scheme; it does
not implement domain decomposition across multiple devices, shared-memory kernel
tiling \cite{micikevicius2009}, or adaptive time stepping. Nor does it implement an
implicit scheme such as Crank--Nicolson \cite{crank1947}, which removes the CFL
restriction at the cost of solving a linear system every step; the explicit scheme
used here trades that restriction for a kernel with no inter-cell solve, which suits
the goal of a straightforward CPU/GPU comparison. The verification strategy, checking
both backends independently against a closed-form solution rather than diffing them
against each other, depends on having a problem with a known analytic solution; it
would not directly generalize to a production case with complex geometry or material
heterogeneity without a different reference (e.g. a manufactured solution or a
trusted reference implementation).

\section{Conclusion}

Verifying both backends independently, rather than trusting a speedup number on its
own, paid off twice in this work: it caught a real GPU boundary-kernel defect
(Section~\ref{sec:bug}), and it caught a flawed convergence-study design before that
flaw could produce a misleadingly clean "order $\approx 4$" result instead of the
genuine order $\approx 2$ (Section~\ref{sec:convergence}). With those issues resolved,
the GPU backend passes all verification checks -- the analytic-solution, conservation,
boundary-invariant, and convergence-order checks alike -- and delivers a
$65$--$68\times$ speedup over the OpenMP/SIMD CPU backend at grid sizes of $128^3$ and
above, a result we can now attribute analytically to the kernel's low arithmetic
intensity ($\approx 0.5$ FLOP/byte) rather than assert without support.

\begin{thebibliography}{9}

\bibitem{micikevicius2009}
P. Micikevicius, ``3D finite difference computation on GPUs using CUDA,'' in
\emph{Proceedings of the 2nd Workshop on General Purpose Processing on Graphics
Processing Units (GPGPU-2)}, 2009.

\bibitem{courant1928}
R. Courant, K. Friedrichs, and H. Lewy, ``\"Uber die partiellen Differenzengleichungen
der mathematischen Physik,'' \emph{Mathematische Annalen}, vol. 100, no. 1, 1928.

\bibitem{obrien1951}
G. G. O'Brien, M. A. Hyman, and S. Kaplan, ``A study of the numerical solution of
partial differential equations,'' \emph{Journal of Mathematics and Physics}, vol. 29,
no. 1, 1951.

\bibitem{nvidia-cuda}
NVIDIA Corporation, \emph{CUDA C++ Programming Guide}, 2024.

\bibitem{dagum1998}
L. Dagum and R. Menon, ``OpenMP: an industry standard API for shared-memory
programming,'' \emph{IEEE Computational Science and Engineering}, vol. 5, no. 1, 1998.

\bibitem{oberkampf2002}
W. L. Oberkampf and T. G. Trucano, ``Verification and validation in computational
fluid dynamics,'' \emph{Progress in Aerospace Sciences}, vol. 38, no. 3, 2002.

\bibitem{williams2009}
S. Williams, A. Waterman, and D. Patterson, ``Roofline: an insightful visual
performance model for multicore architectures,'' \emph{Communications of the ACM},
vol. 52, no. 4, 2009.

\bibitem{crank1947}
J. Crank and P. Nicolson, ``A practical method for numerical evaluation of solutions
of partial differential equations of the heat-conduction type,'' \emph{Mathematical
Proceedings of the Cambridge Philosophical Society}, vol. 43, no. 1, 1947.

\end{thebibliography}

\end{document}
